%======================================================================
%  Guide: rg_session.py (Python, scikit-fem + gmsh)
%======================================================================
\sessiontitle{Guide 3 -- Python: \texttt{rg\_session.py}}
  {Blocks, figures and how to run them. Built on the replica
   \texttt{replica.py} (gmsh + scikit-fem).}

\guidesec{1\quad Running the file}
\begin{enumerate}[itemsep=1pt]
\item Install the packages (any Python $\ge3.9$):
\begin{lstlisting}[style=shstyle]
pip install numpy scipy matplotlib scikit-fem gmsh
\end{lstlisting}
On Linux, the \texttt{gmsh} wheel needs the system libraries of OpenGL
(\texttt{libglu1-mesa}, \texttt{libxcursor1}, \texttt{libxinerama1},
\texttt{libxft2}).
\item Run it, with the noise level and the folder of the figures as optional
arguments:
\begin{lstlisting}[style=shstyle]
python3 rg_session.py              # eps = 0, figures in .
python3 rg_session.py 0.01 figs    # eps = 0.01, figures in figs/
RG_NELEM=60 python3 rg_session.py  # finer mesh
\end{lstlisting}
\item The figures are written to \texttt{fig1.pdf} \dots\ \texttt{fig9.pdf}; the
messages \texttt{FIG k : ...} are printed when each is saved. The run takes
about three seconds.
\end{enumerate}

\guidesec{2\quad Parameters (Block 0)}
\begin{center}\small
\begin{tabular}{>{\ttfamily}l l p{8cm}}
\toprule
\rmfamily name & default & meaning\\
\midrule
ldom, hdom & 1.5, 2.0 & width and height of the rectangle\\
nelem & 30 & mesh size $\ell/\texttt{nelem}$, $\ell_\Gamma/(\texttt{nelem}/2)$ on the crack (\texttt{RG\_NELEM})\\
A, B & crack tips & the true crack (mesh and errors only)\\
eps & 0 & relative noise on the measured values (first argument)\\
seed & 1 & seed of the noise: change it to draw again\\
nxi, xiest & 8, 2.0 & Fourier fields $\xi=1\dots\texttt{nxi}$; $\xi$ of the length estimate\\
nsine & 5 & number of sine--sinh terms\\
nbeta & 12 & turning potential: angles $0,15,\dots,165$ degrees\\
outdir & . & folder of the figures (second argument)\\
\bottomrule
\end{tabular}
\end{center}

\guidesec{3\quad The blocks}
\block{Block 1 -- Geometry and mesh.} The gmsh model of \texttt{replica.py}
(and of the FEniCS file): rectangle, embedded crack, \texttt{Crack} plugin.
The nodes and triangles are read back from gmsh and give a scikit-fem
\texttt{MeshTri}; nodes not used by a triangle are dropped.

\block{Block 2 -- Model.} \texttt{Basis(mesh, ElementTriP2())} and the
stiffness matrix of $\nabla u\cdot\nabla v$, without boundary conditions (nodal
fluxes $q=\mat Ku$). \texttt{basis.doflocs} gives the node coordinates;
\texttt{upnode} the nodes of the cells on the side of $\vn$. FIG~1.

\block{Block 3 -- Procedures.} \texttt{rgap(ub, qb, v)} is
$\sum_{\partial\Omega}u(\mat Kv)-qv$; \texttt{measure} adds the noise to the
measured values; \texttt{opening(u)} pairs the doubled nodes of the crack by
their abscissa and returns $u^+-u^-$; \texttt{dirichlet(mask, values)} solves
with \texttt{condense} and \texttt{solve}.

\block{Block 4 -- Experiment E1.} \texttt{u1 = dirichlet(bot | top, y)}. FIG~2
and FIG~3 (opening, its integral, the infinite-body curve).

\block{Block 5 -- Measurement.} $u$ on the vertical sides without corners,
$q$ on the electrodes; $\RG(1)$.

\block{Blocks 6--7.} Normal, line, moments, lengths, tips; FIG~4 and FIG~5.

\block{Block 8.} (a) Fourier fields, $C/m_0$ and $S/m_0$ against the true
$C/m_0$ of the opening on the lips, length from $\xi=\texttt{xiest}$; FIG~6.
(b) Sine--sinh series as in Cast3M; FIG~7.

\block{Block 9.} Two solves \texttt{ux}, \texttt{uy}; for each angle the noisy
flux, the row of $R$; \texttt{numpy.linalg.svd(R)}, the normal, the weights,
the virtual experiment; FIG~8 and FIG~9.

\guidesec{4\quad The figures}
Those of the common part (Section 4), which were produced by this file:
\texttt{triplot} (FIG~1), \texttt{tricontourf} (FIG~2), line plots (FIG~3, 5--9),
the geometry (FIG~4).

\guidesec{5\quad The printed messages}
One line per block, \texttt{BLOCK k ...}, with the values of Section 5 of the
common part. Signs: $\vn$ is oriented so that $m_0>0$. With
\texttt{eps = 0.01} and \texttt{seed = 1}, for instance: angle error
$4.36^\circ$ (E1) against $0.64^\circ$ (the twelve turning experiments), dipole
length $0.331$, second moment failed ($\mu_2<0$).

\guidesec{6\quad Python pitfalls}
\begin{itemize}[itemsep=1pt]
\item \textbf{Nodes of P2.} \texttt{basis.doflocs} lists the vertices first,
then the midpoints of the edges; the colour plots use only the vertices.
\item \textbf{The lips.} After the \texttt{Crack} plugin, two nodes have the same
coordinates; any search by coordinates must keep both (\texttt{opening}) and
tell them apart by the cells they belong to (\texttt{upnode}).
\item \textbf{Unused gmsh nodes.} gmsh keeps the nodes of the
one-dimensional crack elements; they belong to no triangle and must be
dropped, or the matrix is singular.
\item \textbf{Noise and seed.} The generator is created once
(\texttt{default\_rng(seed)}): Blocks 5 and 9 draw from it in turn.
\end{itemize}

\guidesec{7\quad Listing of \texttt{rg\_session.py}}
\lstinputlisting[style=pystyle]{python/examples/ch3/rg_session.py}
